\chapter{Description of EIRENE input file} \label{sec2}

\textbf{General Remarks} 

\medskip 

Reading of all input data is initialized in subroutine  INPUT. Starting with EIRENE 
version EIRENE\_JSON\_input two I/O options are available: the traditional 
formatted input or the more modern JSON input. In order to distinguish which
type of input is available the code reads the first line of input from I/O channel 
IUNIN. Depending on the first character found the appropriate read procedure is 
carried out. \\
First character 
\begin{description} 
  \itemsep0em
  \item [$= *$] means traditional formatted input is used.
  \item [$= \{$] means JSON input is expected. 
\end{description}
Independent of the input format used the code will produce a corresponding
JSON file. Thus the code works as an input converter as it is planned to completely
move to JSON format over time.

\medskip

\textbf{Formatted Input}

\medskip

The formatted input is read from unit IUNIN via subroutine READ\_FIXFORM

There are 14 "blocks", each one starting with a card 

\begin{lstlisting}
*** n. TEXT
\end{lstlisting} 
where n is the number of the input block, n = 1, 2, \dots, 14.
Some blocks are subdivided into so called "sub-blocks", each one starting with
a card 

\begin{lstlisting}
** TEXT .
\end{lstlisting} 

Each input block is described in one of the following 14
\crefrange{sec2.1}{sec2.14}:

\medskip

\textsl{2.1  Input data for operating mode}

\textsl{2.2  Input data for standard mesh}

\textsl{2.3A Input data for "Non-default Standard Surfaces"}

\textsl{2.3B Input data for "Additional Surfaces"}

\textsl{2.4  Input data for species specification and atomic physics module}

\textsl{2.5  Input data for plasma background}

\textsl{2.6  Input data for surface interaction models}

\textsl{2.7  Input data for initial distribution of test particles}

\textsl{2.8  Additional data for some specific mesh zones}

\textsl{2.9  Data for statistics and non-analog methods}

\textsl{2.10 Data for additional volumetric and surface averaged tallies}

\textsl{2.11 Data for numerical and graphical output}

\textsl{2.12 Data for plasma diagnostic module DIAGNO}

\textsl{2.13 Data for nonlinear and time-dependent mode}

\textsl{2.14 Data for interfacing with external codes}

\medskip

This last block 14 may be read in from the user supplied part (with any format
specified there) from this same unit IUNIN, (but equally well from any other
unit specified there), e.g., from the subroutine INFCOP for interfacing with
other codes.

\medskip 

\textbf{Input format} 

\medskip 

The following conventions are used:
\begin{itemize}
  \item Standard Fortran naming conventions, i.e., variable names starting with
        the letters I, J, \dots, N are Integer, all other variables are Real.
\end{itemize}
In addition to these standard convention the following rules are employed:
\begin{itemize}
  \item Character strings start with the letter C\dots\ or with T\dots
  \item
        Logical variables start with NL\dots, LG\dots\ or with TRC\dots
\end{itemize}
The format for a card containing only real variables: (6E12.4)

The format for a card containing only integer variables: (12I6)

The format for a card containing only logical variables: (12(5L1,1X))

The format for a card containing text: (A72) 

The format for a card containing K (K $<$ 6) integer variables first and then
up to 6-K real variables next: (K(I6,6x),(6-K)E12.4) 

An input card starting with '*' is recognized as a comment line, i.e., it does
not contain input data.
These comment lines should not start with  '**' or '***', because in this
latter case EIRENE would assume that the input for the current block is
completed.
It would expect the first card of the following block (or '* comment\dots'
lines belonging to this next block) as next card.
Any data not read from the input file, before the following block is read, are
set to default values.

An arbitrary number of such comment lines may be included in the input file,
but each block of comment lines must always be put at the beginning of a block
or sub-block.
I.e., comment lines are only identified by EIRENE if preceded by a card
starting with '**' or '***' indicating a new sub-block or block respectively.

\medskip

Example (see \cref{sec2.7} below):

\begin{lstlisting}
*** 7. Data for initial distribution of test particles
* this is a comment
* there are NSTRAI=2 strata in this run
* this is a comment
 2
** Data for first stratum
* this is a surface recycling source of helium ions, IPLS=2
* this is a comment
FFFF
  .
  .
  .
  .
** Data for second stratum
* this is a hydrogen ion volume re-combination source, IPLS=1
FFFF
  .
  .
  .
*** 8. Data for some specific zones
  .
  .
  .
  .
\end{lstlisting}

and so on.

\medskip

In very few cases formats other than these are used.
They are then explicitly mentioned in this manual and in case of doubt the user
should check in subroutine READ\_FIXFORM.

\medskip

\medskip

\textbf{JSON input}

\medskip

JSON is an acronym for JavaScript Object Notation \cite{kn:Json} which is an open standard 
file format.  
It can represent four primitive types and  two structured types. The primitive types are
strings, numbers, boolean and null. The structured types are objects and arrays. 

An object is a

\begin{center}
name : value
\end{center}

pair where "name" is a string and "value" can be any of the primitive or structured types.
A JSON input is easily readable and can also be viewed within browsers or online editors.

Libraries for JSON have been implemented for many programming languages including
FORTRAN. An implementation for an object-oriented FORTRAN library can be found on github 
\cite{kn:Williams}.
 
Unlike the formatted I/O reading the JSON input file does not use the I/O unit IUNIN but 
expects the input file to have a specific filename. The name \textbf{eirene.input.json} has
been chosen. Therefore this file should be linked to unit IUNIN initially.

Similar to the formatted input the JSON input consists of 15 blocks. A zeroth block was
introduced to hold the input header and comments on the specified case. Each block is its own 
JSON object. Usually all blocks are placed into the main input file \textbf{eirene.input.json} 
but can be put into separat files as well. Then the main input file only gets a skeleton object of 
the block with an indicator where to look for the corresponding input. 
The names of the blocks follow the naming in the formatted input and describe the intended
subject of the data in the block. An entry "MANUAL" in each block refers to the corresponding
section in this manual. Sub-blocks are grouped into their own objects with appropriate names.
The blocks are named:

\medskip

\textsl{HEADER}

\medskip

\textsl{GENERAL\_DATA}

\medskip

\textsl{STANDARD\_MESH} \\
	\hspace*{1cm} \textsl{RADIAL\_GRID} \\
	\hspace*{1cm} \textsl{POLOIDAL\_GRID} \\
	\hspace*{1cm} \textsl{TOROIDAL\_GRID} \\
	\hspace*{1cm} \textsl{MESH\_MULTIPLICATION} \\
	\hspace*{1cm} \textsl{ADD\_CELLS}

\medskip

\textsl{NONDEF\_STD\_SURFACES}

\medskip

\textsl{ADDITIONAL\_SURFACES}

\medskip

\textsl{PHYSICS\_MODEL} \\
	\hspace*{1cm} \textsl{SPECIES\_SPEC} \\
   		\hspace*{2cm} \textsl{REACTIONS} \\
		\hspace*{2cm} \textsl{ATOMS} \\
		\hspace*{2cm} \textsl{MOLECULES} \\
		\hspace*{2cm} \textsl{TEST\_IONS} \\
		\hspace*{2cm} \textsl{PHOTONS} \\
	\hspace*{1cm} \textsl{BACKGROUND} \\
		\hspace*{2cm} \textsl{BULK\_IONS} \\
		\hspace*{2cm} \textsl{PLASMA} 

\medskip

\textsl{REFLECTION\_MODELS}

\medskip

\textsl{SOURCES}

\medskip

\textsl{SPECIFIC\_ZONES}

\medskip

\textsl{STATISTICS}

\medskip

\textsl{ADDITIONAL\_TALLIES} 

\medskip

\textsl{OUTPUT}

\medskip

\textsl{DIAGNOSTICS}

\medskip

\textsl{TIMEDEPENDENT\_MODE}

\medskip

\textsl{INTERFACING}

\medskip

The input of interfacing blocks has been converted to JSON format as well. Any other data
read by user routines is placed into a file \textbf{user\_data.input} by the converter and is read
in the same way as with the formatted input.

The names for the variables are almost everywhere identical to those given in this manual. 
Deviations are indicated in the corresponding block.

Variables and blocks can be omitted from the input file if they are not required by the 
simulation. Omitted variables are set to 0 or .false. in case of boolean. The converter however
will always produce the complete JSON files. The two blocks \textsl{ADDITIONAL\_SURFACES} 
and \textsl{PHYSICS\_MODEL} are always placed into and read from separate files named 
\textbf{eirene\_add\_surfaces.input.json} and \textbf{eirene\_physics\_model.input.json}.


\medskip

\medskip

\textbf{Units} 

cgs units are used, with two exceptions:

\medskip

All \textbf{temperatures and energies} are in \si{\electronvolt}.

All \textbf{particle fluxes} are in Ampere (i.e.\
\SI{1.6022e-19}{\#\per\second}), even neutral particle fluxes.
Since energies are given in \si{\electronvolt} units, a particle flux $\times$
energy is already in \si{\watt}.

E.g., to convert a \textbf{gas flow rate} $Q$ (= ``mass-flow rate" at
temperature $T_0$), which is a quantity given with dimension pressure times
volume divided by time, to a particle flux (source strength, given in
1/\;time), one utilizes the universal gas law $P \times V = N \times k_B T$ to
convert from pressure times volume (typically at standard temperature $T_0=
  \SI{273.15}{\kelvin}$) to the number $N$ of particles.
For example  $x$ \ac{SLPM} first divide $x$ by \num{22.4} (then:
\si{\mole\per\minute}, at ``standard conditions" = normal pressure
\SI{1013.25}{\milli\bar} and normal temperature $T_0$).
Next multiply by Avogadro's no.\ $N_A = \num{6.0221e23}$, (i.e., then
particles~\si{\per\minute}), then divide by \num{60} (then:
particles~\si{\per\second}), and finally multiply by the elementary charge
\num{1.6022e-19}
(then \si{\ampere}, as used for particle fluxes in EIRENE).
Gas (mass-) flow rates given in other units, e.g.\
\si{\torr\cubic\metre\per\second} are converted correspondingly.

Another frequently used dimensional quantity is the (effective) 
\textbf{pumping speed} $S$, and also \textbf{conductance} $L$ which both have 
dimension volume / time (e.g.\ litres/sec), i.e.\ these are 
``volumetric flow rates" rather than mass flow rates.
Multiplying these volume flow rates by the entrance gas pressure (or by the
pressure difference in case of conductance) produces the (pump-) throughput, a
mass flow rate again, i.e., at given temperature a ``pumped particle flux"
(\si{\per\second}).

Pumping speeds and conductances typically depend on pressure themselves, and on
temperature $T$ and the type of gas (mass $m$).
In the molecular flow regime the pressure dependence in $S, L$ diminishes, and
$S$ ($L$) can be written as product $S,L =A_{S,L} \cdot v$, of an area and
velocity, with $v \propto \sqrt{T/m}$.
Implementation of a specified volumetric flow (e.g.\ a given effective pumping
speed at pumps) in EIRENE is via effective surface area and  sticking
probabilities, see \cref{sec2.6.1}, input flag \textbf{RECYCT}

All \textbf{densities} are in \si{\per\cubic\centi\metre}.

All \textbf{velocities} are in either \si{\centi\metre\per\second} or in
(isothermal) Mach-number units.

All geometrical data are in \si{\centi\metre}.

Note: consequently, e.g., energy densities are in
$\si{\electronvolt\per\cubic\centi\metre}$, and energy fluxes are given in
\si{\electronvolt\ampere}, i.e., in \si{\watt}.

Hence: In order to convert from EIRENE \textbf{energy density} units
(\si{\electronvolt\per\cubic\centi\metre}) into  \textbf{pressure} units in
\si{\newton\per\square\metre} (i.e., Pascal: \si{\pascal}) one must multiply by
\num{1.6022e-19} (\si{\electronvolt} to \si{\newton\metre}) and multiply by
\num{1e6} ($\si{\per\cubic\centi\metre} \rightarrow \si{\per\cubic\metre}$).
Check: \SI{1}{\pascal} gas at \SI{300}{\kelvin} corresponds to about
$\num{2.7e14} \ \times \ \SI{0.026}{\electronvolt\per\cubic\centi\metre}
  \approx \SI{7e12}{\electronvolt\per\cubic\centi\metre}$, or:
a gas pressure P in \si{\pascal}, at gas temperature T (\si{\kelvin})
corresponds to a gas density n with $n (\si{\per\cubic\metre}) = P/k_B T$,
with $k_B = \num{1.38e-23}
$ the Boltzmann constant in \si{\joule\per\kelvin} units.
At \SI{300}{\kelvin} and \SI{1}{\pascal}, the corresponding gas density is
%$n=\SI{2.415 \dots e20}{\per\cubic\metre} = \SI{2.415 \dots e14}{\per\cubic\centi\metre}$.
$n=\num{2.415} \dots \SI {e20}{\per\cubic\metre} = \num{2.415} \dots \SI{e14}{\per\cubic\centi\metre}$.

For example, energy densities as computed by EIRENE defaults tallies 5, 6, 7
and 8 (see \cref{tab2}) must be converted into ``temperature densities" (a
factor 2/3, if one neglects the energy of the directed motion) and then be
re-scaled to pressure units as described above.

For further scaling to other pressure units note, e.g.:

\SI{101.3}{\pascal} = \SI{760}{\milli\torr} = \SI{1}{\milli\bar}

Further: momentum is in \si{\gram\centi\metre\per\second}, and a
\textbf{momentum flux} or momentum rate of change is in
\si{\gram\centi\metre\per\second\ampere} (per cell) and the momentum source
density in the EIRENE tallies for interfacing with plasma codes is then in
\si{\gram\centi\metre\per\second\ampere\per\cubic\centi\metre}.

Hence, conversion from \si{\gram\centi\metre\per\second\ampere} into
\si{\pascal\square\metre} (momentum source per cell in plasma momentum balance
equations in SI units) is done by firstly multiplying with \num{1e-5} (then:
\si{\kilo\gram\metre\per\second\ampere}) and then dividing by \num{1.6022e-19}
(hence: \si{\kilo\gram\metre\per\square\second}).
After dividing by the cell volume (\si{\cubic\metre}) the momentum source is in
\si{\pascal\per\metre}, or \si{\newton\per\cubic\metre} (force density).

\medskip

